% Generated by report/tools/tables.py from the result files of the repository. Do not edit.
\begin{table}[tbp]
\centering
\footnotesize
\setlength{\tabcolsep}{5pt}
\caption[La règle « tranché » sur des hypothèses nulles]{La règle « tranché » mise à l'épreuve sur de vraies hypothèses nulles : deux groupes disjoints du même programme, où aucune différence n'existe. Sur chaque tirage, intervalle bootstrap percentile de la probabilité d'amélioration et intervalles des deux IQM (2\,000 rééchantillons, chaque groupe rééchantillonné à part, méthode de rliable réécrite en NumPy, \texttt{default\_rng(2026)}). Colonnes : part des tirages où l'intervalle de la probabilité d'amélioration exclut 0,5, où les intervalles des IQM sont disjoints, où l'un ou l'autre se produit (la règle du protocole, \textbf{en gras}). Une règle exacte au niveau de 95~\% se tromperait dans 5~\% des tirages : la règle est un peu libérale à 10 contre 10, nettement à 5 contre 5. L'excès vient de la condition sur la probabilité d'amélioration. Source : \file{sim/analysis/rule\_calibration.json} (\file{sim/analysis/rule\_calibration.py}).}
\label{tab:regle-tranche}
\begin{tabular}{@{}lrrrr@{}}
\toprule
Groupes du même programme & \makecell[r]{IC de $P$\\{}exclut 0,5} & \makecell[r]{IC des IQM\\{}disjoints} & \makecell[r]{L'un ou l'autre\\{}(la règle)} &  \\
\midrule
\multicolumn{5}{@{}l}{\emph{Les 30 essais de HalfCheetah (E1 et B1), 2\,000 tirages}} \\
\quad PPO PyTorch, 10 contre 10 & 5,4~\% & 0,6~\% & \textbf{5,4~\%} &  \\
\quad PPO PyTorch, 5 contre 5 & 9,9~\% & 2,6~\% & \textbf{9,9~\%} &  \\
\quad TRPO, 10 contre 10 & 6,2~\% & 0,9~\% & \textbf{6,2~\%} &  \\
\quad TRPO, 5 contre 5 & 9,1~\% & 2,2~\% & \textbf{9,1~\%} &  \\
\midrule
\multicolumn{5}{@{}l}{\emph{La règle à 5 contre 5 sur les seize couples de E2 (graines 0 à 19), 2\,000 tirages}} \\
Programme & HalfCheetah & Hopper & Walker2d & Swimmer \\
\cmidrule(lr){1-5}
\quad PPO PyTorch & 10,1~\% & 9,2~\% & 9,3~\% & 9,3~\% \\
\quad PPO JAX & 8,6~\% & 9,5~\% & 9,6~\% & 9,8~\% \\
\quad SAC & 9,9~\% & 9,9~\% & 8,6~\% & 8,6~\% \\
\quad CrossQ & 9,6~\% & 9,0~\% & 9,9~\% & 10,1~\% \\
\multicolumn{5}{@{}l}{Médiane 9,5~\% (8,6~\% à 10,1~\%)} \\
\bottomrule
\end{tabular}
\end{table}
